\section*{Chapter \ref{ch:g6:reading-economic-charts} --- Reading Economic Charts}

\begin{solution}{exo:g6:reading-economic-charts:1}
(a) About $1.5$ billion people in 1900 and about $6$ billion in 2000.
(b) About $6 - 1.5 = 4.5$ billion more people.
\end{solution}

\begin{solution}{exo:g6:reading-economic-charts:2}
(a) The second village: $200\,000$ kg against $120\,000$ kg.
(b) First village: $120\,000 \div 400 = 300$ kg per person; second village:
$200\,000 \div 1\,000 = 200$ kg per person. The first village, with the smaller
total, has the larger harvest per person.
\end{solution}

\begin{solution}{exo:g6:reading-economic-charts:3}
(a) East Asia in 1820 ($41.1\,\%$); South and South-East Asia in 2022
($32.6\,\%$). (b) $41.1 - 21.0 = 20.1$ percentage points.
\end{solution}

\begin{solution}{exo:g6:reading-economic-charts:4}
(a) Walking: $\frac{15}{25} \times 100\,\% = 60\,\%$; bus:
$\frac{6}{25} \times 100\,\% = 24\,\%$; car: $\frac{4}{25} \times 100\,\% = 16\,\%$.
(b) $60 + 24 + 16 = 100$: the three ways cover every pupil once.
\end{solution}

\begin{solution}{exo:g6:reading-economic-charts:5}
(a) The change is $31\,946 - 11\,061 = 20\,885$ dollars, and
$20\,885 \div 11\,061 \approx 1.89$: a rise of about $189\,\%$.
(b) $31\,946 \div 11\,061 \approx 2.9$: income per person was almost three times
as large. (A rise of $189\,\%$ and ``$2.9$ times as large'' say the same thing:
$100\,\% + 189\,\% = 289\,\%$ of the 1950 value.)
\end{solution}

\begin{solution}{exo:g6:reading-economic-charts:6}
(a) Between 1\,000 and 10\,000 (about 3\,000). (b) Between 1\,000 and 10\,000
(about 7\,500). (c) Between 10\,000 and 100\,000 (just under 40\,000).
(d) Around 1940.
\end{solution}

\begin{solution}{exo:g6:reading-economic-charts:7}
(a) Each gridline is ten times the one below: two gridlines up is
$10 \times 10 = 100$ times larger. (b) Two: from 1\,000 to 10\,000 and from
10\,000 to 100\,000. (c) $1\,000 = 10 \times 10 \times 10$: three gridlines.
\end{solution}

\begin{solution}{exo:g6:reading-economic-charts:8}
(a) That the number of people in the world rose, from about 1 billion in 1820
to almost 8 billion in 2022. (b) The chart is about the number of people only;
it contains no price. More people could live with food that became cheaper,
or with food that became dearer, or with food whose price did not change.
(c) A chart of the price of food over the same years, compared with what
people earn: for example the number of hours of work needed to buy a kilogram
of bread or of rice.
\end{solution}

\begin{solution}{exo:g6:reading-economic-charts:9}
(a) Northland's bar rises $9\,500 - 9\,000 = 500$ dollars above the bottom of
the axis, Southland's $10\,500 - 9\,000 = 1\,500$ dollars: Southland's bar looks
$1\,500 \div 500 = 3$ times as tall. (b) $10\,500 \div 9\,500 \approx 1.11$:
Southland's income per person is only about $1.11$ times Northland's (about
$11\,\%$ higher). (c) Where the numbers come from (there is no \omterm{def:g6:reading-economic-charts:source-line}{source line});
why the axis starts at 9\,000; in what kind of dollars the incomes are
measured, and for which year.
\end{solution}

\begin{solution}{exo:g6:reading-economic-charts:10}
(a) Island A has the larger total income ($30\,000$ against $10\,000$ million
dalens). Per person: A has $30\,000\,000\,000 \div 2\,000\,000 = 15\,000$ dalens,
B has $10\,000\,000\,000 \div 500\,000 = 20\,000$ dalens: B has the larger
income per person. (b) $2\,200\,000$ people and $33\,000$ million dalens:
$33\,000\,000\,000 \div 2\,200\,000 = 15\,000$ dalens, unchanged. (c) Total
income would have to rise by more than $10\,\%$, faster than the population.
\end{solution}

\begin{solution}{exo:g6:reading-economic-charts:11}
(a) Yes: $3\,000$ is ten times $300$ and $30\,000$ is ten times $3\,000$, and on a
\omterm{def:g6:reading-economic-charts:times-ten-scale}{times-ten scale} equal ratios are equal distances. (b) Neither: both steps
double the value, so they have the same length. (c) On an ordinary scale, from
$5\,000$ to $10\,000$ (a step of $5\,000$) is five times as long as from $1\,000$ to
$2\,000$ (a step of $1\,000$). (d) The quantity is multiplied by the same number
in each equal stretch of time.
\end{solution}

\begin{solution}{exo:g6:reading-economic-charts:12}
(a) 1820: $1.042 \times \frac{13.4}{100} \approx 0.14$ billion people;
2022: $7.802 \times \frac{5.5}{100} \approx 0.43$ billion people. (b) No: the
number of people in Western Europe was multiplied by about three; only its
\omterm{def:g6:reading-economic-charts:share}{share} fell, because the rest of the world grew faster. (c) Its \omterm{def:g6:reading-economic-charts:share}{share} would have
to fall \emph{and} its number of people would have to fall. The \omterm{def:g6:reading-economic-charts:share}{shares} chart
alone cannot show it: a \omterm{def:g6:reading-economic-charts:share}{share} can fall while the number rises, so one also needs
the total (here, the world population chart) or a chart of the numbers
themselves.
\end{solution}

\begin{solution}{pb:g6:reading-economic-charts:1}
\textbf{1.} $18\,000 + 12\,000 + 6\,000 + 24\,000 = 60\,000$ people.
\textbf{2.} Harbour $\frac{18\,000}{60\,000} = 30\,\%$; Old Town $20\,\%$;
Hillside $10\,\%$; Riverside $40\,\%$.
\textbf{3.} $30 + 20 + 10 + 40 = 100$. The bar has four pieces of heights 30, 20,
10 and 40 out of 100.
\textbf{4.} Harbour $360\,000\,000 \div 18\,000 = 20\,000$ dalens; Old Town
$300\,000\,000 \div 12\,000 = 25\,000$; Hillside $240\,000\,000 \div 6\,000 =
40\,000$; Riverside $360\,000\,000 \div 24\,000 = 15\,000$ dalens.
\textbf{5.} $360 + 300 + 240 + 360 = 1\,260$ million dalens;
$1\,260\,000\,000 \div 60\,000 = 21\,000$ dalens per person.
\textbf{6.} $\frac{360}{1\,260} \times 100\,\% \approx 28.6\,\%$ of the income,
for $40\,\%$ of the people.
\textbf{7.} By total income: Harbour and Riverside (360 each), then Old Town
(300), then Hillside (240). By income per person: Hillside (40\,000), Old Town
(25\,000), Harbour (20\,000), Riverside (15\,000). Hillside goes from last to
first: it has few people, each with a high income. Riverside goes from first
place (shared) to last: a large total shared among many people.
\textbf{8.} Right: $21\,000$ dalens is the town's income per person, its
average. Wrong: it is not what ``everyone'' earns. The districts range from
$15\,000$ to $40\,000$ dalens per person, and inside each district incomes differ
again; an average does not show how a total is shared.
\textbf{9.} $\frac{60\,000 - 6\,000}{6\,000} \times 100\,\% = 900\,\%$.
\textbf{10.} $60\,000 \div 600 = 100$ times.
\textbf{11.} $600$ lies between 100 and 1\,000; $6\,000$ between 1\,000 and
10\,000; $60\,000$ between 10\,000 and 100\,000, each at the same place within
its step. The points are evenly spaced, one step apart: in each fifty-year
period the population was multiplied by the same number, ten.
\textbf{12.} $600 \div 60\,000 = 0.01$: the 1925 point is at $1\,\%$ of the
height, almost on the bottom of the axis, and the 1975 point at $10\,\%$. A
reader would think that the town hardly grew between 1925 and 1975, when in
fact it was multiplied by ten, exactly as in the next fifty years.
\end{solution}
